\chapter{Nota da autora}

\noindent{}Pegue o primeiro texto que encontrar: \emph{texto-memória},
\emph{texto-qualquer}, \emph{texto-nenhum}. Ele sempre será um mapa. Foi
o que me ocorreu ao ler o retrato de adeus e de amor de minha trisavó
Dvora\footnote{Dvora é mãe de Aba, que é pai de Sarah Dora, que é mãe de
  Luciana, que é mãe de Tamara --- que nasceu na cidade de São Paulo, em
  1999.}. Eu o acolhi em mim como uma espécie de guia, uma casa, um
mundo por completo.

Agora é hora de pousar nas páginas, dentro delas. Bicar as palavras um
pouco mais para perto de nós, fazê-las falar o que queremos dizer.
\emph{O que escutamos ao colar o ouvido no espelho?}

\emph{Perspectivas judaico-diaspóricas sobre o lugar} investiga a
maneira pela qual a cultura e as tradições judaicas se apropriam dos
espaços de trânsito, os movimentos cíclicos\footnote{Por meio da
  ``consagração'' de um lugar há a repetição da cosmogonia, ou seja, a
  vitória do \emph{lugar sagrado} sobre o \emph{espaço profano} deve ser
  repetida simbolicamente todos os anos, pois todos os anos o mundo
  necessita ser criado novamente desde o seu princípio.} e constantes
das diásporas, e os tornam uma possibilidade de pertencimento --- um
começo, um meio e um começo.

No espaço \emph{talmúdico}\footnote{A organização espacial da página,
  nos textos judaicos religiosos, é feita de tal maneira que as diversas
  interpretações rabínicas são dispostas sem hierarquização. Assim, no
  \emph{Talmud}, todos os intérpretes são convidados a participar
  ativamente do texto, em um esforço para esgotar o sentido e as
  ramificações deste material.} deste texto há a intenção de definir
lugar, ainda que na ausência de um. É necessariamente neste ponto --- em
que não se sabe ao certo o que se encontrará --- que imaginamos ter algo
a dizer. Pois, como intérprete de uma longa tradição de textos e de
livros, acredito que o nosso mundo --- e todas as coisas que ele contém
--- está sendo criado e reinventado por nós a todo instante.
